\chapter{Experimental Results}

\section{TabFact Data Analysis}

As is standard practice in text classification using a dataset, analysing the dataset to get insight to confirm the dataset you are using is suitable for a general application of this problem and to give you insight for training or fine-tuning is necessary. The original paper already displayed a good analysis of the data. I recreated their analysis in tables Table~\ref{tab:tabfact_train_val_test_split} and Table~\ref{tab:tabfact_simple_complex_split}. I also performed additional analysis to confirm the strength of this dataset as a suitable benchmark. My recreated analysis showed differences in row and column count, but as the number of data being read in was correct, this can be attributed to the text-flattening process of the tables in the JSON from the original HTML. My data on the simple and complex channels show more accurate statistics than in the original paper because they filtered out 400 tables and did not recalculate the metrics. 
\begin{table}[ht]
\center
\begin{tabular}{lcccc}
\hline
Split & \#Sentence & \#Table & Rows & Columns\\
\hline
Train & 92,283 & 13,182 & 14.52 & 6.35\\
Val & 12,792 & 1,696 & 14.70 & 6.24\\
Test & 12,779 & 1,695 & 14.43 & 6.25\\
\hline
\end{tabular}
\caption{Sentences and tables in each dataset.}
\label{tab:tabfact_train_val_test_split}
\end{table}
\begin{table}[ht]
\center
\begin{tabular}{lcccc}
\hline
Channel & \#Sentence & \#Table & Len(Ent) & Len(Ref)\\
\hline
Simple (r1) & 49,850 & 9,181 & 13.05 & 13.32\\
Complex (r2) & 68,004 & 7,392 & 14.19 & 14.15\\
Total & 117,854 & 16,573 & 13.71 & 13.79\\
\hline
\end{tabular}
\caption{Simple and complex channel sentences and tables.}
\label{tab:tabfact_simple_complex_split}
\end{table}

\subsection{Entailed and Refuted Statement Distribution across Simple and Complex Channels}
I performed a more detailed breakdown of the distribution of entailed statements as well as that of the r1 (simple channel) and r2 (complex channel) that TabFact organised its data into, displayed in Table~\ref{tab:tabfact_ent_ref_breakdown}. As you can see from the table, there is a slight imbalance in the number of entailed and refuted statements in favour of entailed in the training set. The validation and test datasets are nearly completely balanced. There are more complex channel statements overall; these are statements which require complex understanding of the table structure, such as arithmetic operations on a column to get a conclusion. The distribution of complex and simple channel statements is balanced between entailed and refuted statements in the validation and test datasets, but there are 27\% more complex statements that are entailed in the training dataset is the only potential weakness in this dataset, but the rest demonstrates a quality balance. The data in TabFact rewrote positive statements, so it was possibly more difficult to convert some of those r2 entailed statements into refuted ones. 
\begin{table}[ht]
\center
\begin{tabular}{lcc}
\hline
Data & Entailed Statements & Refuted Statements\\
\hline
Train & 50,820 & 41,463 \\
Validation & 6,478 & 6,314 \\
Test & 6,425 & 6,354 \\
Simple channel (r1) & 26,930 & 23,314 \\
Complex channel (r2) & 37,032 & 30,999 \\
Train simple (r1) & 22,528 & 19,114 \\
Train complex (r2) & 28,466 & 22,477 \\
Validation simple (r1) & 2,248 & 2,124 \\
Validation simple (r2) & 4,260 & 4,219 \\
Test simple (r1) & 2,154 & 2,076 \\
Test complex (r2) & 4,306 & 4,303 \\
\hline
\end{tabular}
\caption{Entailed and refuted entries in each data grouping.}
\label{tab:tabfact_ent_ref_breakdown}
\end{table}

\subsection{Using TabFact preprocessing to study average statement complexity}
I took advantage of TabFact’s preprocessing, which included entity linking and part-of-speech labelling in order to perform analysis based on content words and entities in the statements. The entity defines the name of something they need to look up in the table, such as the name of a person, for example. Part-of-speech (POS) labelling  indicates the type of grammatical word is for each word in the sentence for example NN is a singular noun. 

\subsubsection{Entity Analysis}
Displayed in Table~\ref{tab:tabfact_entity_analysis}, entailed and refuted statements have roughly 3.5 entities per statement, so roughly 3-4 elements the fact verification method has to source information for in the table. Most of the entities are 1-2 words in length. TabFact’s entity linking system would identify mostly a cell per entity, across 2-3 unique rows and columns on average. This emphasises the strength of the dataset, as multiple unique elements have to be identified and considered in the entailed or refuted outcome, with compound word entities across several rows and columns showing that fact verification is not trivial and requires a traversal and understanding of the table’s structure.
\begin{table}[ht]
\center
\begin{tabular}{lc}
\hline
Mean Entity Measure & Value \\
\hline
Entities identified per sentence & 3.49 \\
Entities identified per entailed sentence & 3.51 \\
Entities identified per refuted sentence & 3.46 \\
Entity word length & 1.72 \\
Evidence cells identified per sentence & 3.38 \\
Rows identified per sentence & 2.57 \\
Columns identified per sentence & 2.37 \\
\hline
\end{tabular}
\caption{Tabfact entity and evidence retrieval analysis.}
\label{tab:tabfact_entity_analysis}
\end{table}

\subsubsection{POS Tag Analysis}
Part-of-speech tags in tabfact allow you to pick out the content words, i.e. adverbs, adjectives, nouns, in each statement. As you can see in Table~\ref{tab:tabfact_pos_tag_analysis}, the distribution of content words hovers around 7 for all statements across simple and complex channels and entailed and refuted statements. This suggests that complexity for understanding is roughly the same in entailed and refuted statements. There are more complex channel content words, but that is expected, as you will need more content words to command a more complex logical operation, so that is good to see. 
\begin{table}[ht]
\center
\begin{tabular}{lcccc}
\hline
Mean POS Tag Measure & Value\\
\hline
Content words & 7.15\\
Content words in r1 (simple) & 6.50\\
Content words in r2 (complex) & 7.62\\
Content words in refuted statements & 7.23\\
Content words in entailed statements & 7.08\\
\hline
\end{tabular}
\caption{POS Tag analysis across data groupings.}
\label{tab:tabfact_pos_tag_analysis}
\end{table}

\subsubsection{Content Word Distribution}
It was important to analyse the distribution of content words in entailed and refuted statements to ensure there were no obvious words that are skewed in one class that the model can use to figure out if a claim is entailed or refuted instead of using the table to make that conclusion like we want it to. I excluded entities, of course, as this was to study if any surrounding words gave away if it was entailed or refuted. I did this for all words and then targeted the comparative words. Across all words, shown in Table~\ref{tab:tabfact_top_20_content_words} the non-comparative statement distribution is almost identical between entailed and refuted statements, which shows the statements talked about the exact same domains. Counting just the comparative terms, Table~\ref{tab:tabfact_top_20_comparatives} shows that they could not be used to cheat to determine if a claim was entailed or refuted, as there are positive and negative comparative terms in entailed and refuted statements. Therefore, effective fact verification of this dataset must be judged by studying the table and the claim.
\begin{table}[ht]
\centering
\small
\begin{tabular}{r l r l r l r}
\hline
 & \multicolumn{2}{c}{\textbf{All Statements}} & \multicolumn{2}{c}{\textbf{Entailed Statements}} & \multicolumn{2}{c}{\textbf{Refuted Statements}} \\
\cline{2-3} \cline{4-5} \cline{6-7}
\textbf{Rank} & \textbf{Word} & \textbf{Count} & \textbf{Word} & \textbf{Count} & \textbf{Word} & \textbf{Count} \\
\hline
1  & be      & 78,143 & be      & 41,481 & be      & 36,662 \\
2  & have    & 28,405 & have    & 15,391 & have    & 13,014 \\
3  & game    & 15,298 & game    & 8,260  & game    & 7,038  \\
4  & play    & 10,075 & play    & 5,389  & play    & 4,686  \\
5  & only    & 8,451  & only    & 4,523  & more    & 3,951  \\
6  & 2       & 8,051  & 2       & 4,467  & only    & 3,928  \\
7  & 1       & 7,684  & 1       & 4,226  & 2       & 3,584  \\
8  & more    & 7,546  & win     & 3,609  & 1       & 3,458  \\
9  & time    & 6,665  & time    & 3,597  & time    & 3,068  \\
10 & win     & 6,661  & more    & 3,595  & win     & 3,052  \\
11 & team    & 6,543  & team    & 3,504  & team    & 3,039  \\
12 & number  & 5,144  & number  & 2,699  & number  & 2,445  \\
13 & 3       & 5,033  & 3       & 2,667  & 3       & 2,366  \\
14 & highest & 4,799  & highest & 2,597  & highest & 2,202  \\
15 & most    & 3,719  & most    & 2,147  & less    & 1,983  \\
16 & 4       & 3,210  & 4       & 1,731  & most    & 1,572  \\
17 & season  & 2,940  & season  & 1,537  & 4       & 1,479  \\
18 & less    & 2,937  & episode & 1,509  & not     & 1,443  \\
19 & not     & 2,923  & point   & 1,502  & season  & 1,403  \\
20 & point   & 2,794  & score   & 1,498  & lowest  & 1,401  \\
\hline
\end{tabular}
\caption{Top 20 most frequent lemmas across all, entailed, and refuted statements.}
\label{tab:tabfact_top_20_content_words}
\end{table}
\begin{table}[ht]
\centering
\small
\begin{tabular}{r l r l r l r}
\hline
 & \multicolumn{2}{c}{\textbf{All Statements}} & \multicolumn{2}{c}{\textbf{Entailed Statements}} & \multicolumn{2}{c}{\textbf{Refuted Statements}} \\
\cline{2-3} \cline{4-5} \cline{6-7}
\textbf{Rank} & \textbf{Word} & \textbf{Count} & \textbf{Word} & \textbf{Count} & \textbf{Word} & \textbf{Count} \\
\hline
1  & more     & 7,546 & more     & 3,595 & more     & 3,951 \\
2  & highest  & 4,799 & highest  & 2,597 & highest  & 2,202 \\
3  & most     & 3,718 & most     & 2,147 & less     & 1,980 \\
4  & less     & 2,933 & lowest   & 1,245 & most     & 1,571 \\
5  & lowest   & 2,611 & less     & 953   & lowest   & 1,366 \\
6  & least    & 1,265 & higher   & 655   & least    & 687   \\
7  & higher   & 1,249 & least    & 578   & higher   & 594   \\
8  & largest  & 865   & largest  & 463   & largest  & 402   \\
9  & lower    & 714   & lower    & 316   & lower    & 398   \\
10 & greater  & 595   & larger   & 275   & greater  & 395   \\
11 & larger   & 569   & longest  & 249   & fewer    & 364   \\
12 & fewer    & 546   & smallest & 235   & larger   & 294   \\
13 & smallest & 506   & earliest & 229   & smallest & 271   \\
14 & longest  & 420   & greater  & 200   & smaller  & 247   \\
15 & smaller  & 416   & best     & 197   & best     & 198   \\
16 & earliest & 401   & fewer    & 182   & earliest & 172   \\
17 & best     & 395   & smaller  & 169   & longest  & 171   \\
18 & fewest   & 289   & fewest   & 159   & fewest   & 130   \\
19 & latest   & 229   & latest   & 124   & better   & 110   \\
20 & better   & 223   & greatest & 121   & latest   & 105   \\
\hline
\end{tabular}
\caption{Top 20 most frequent comparative and superlative words across all, entailed, and refuted statements.}
\label{tab:tabfact_top_20_comparatives}
\end{table}

\subsection{Caption Analysis}
To again ensure that entailed and refuted claims on their own could not give an indication of the fact verification with the table, I tested semantic similarity with the table topic on cosine similarity between a claim with its table caption. Shown in Table~\ref{tab:tabfact_caption_cosine_sim}, cosine similarity was almost identical for both classes suggesting semantic similarity with the table subject does not give away the verification label. This was performed on the test set to get a quick idea of this. I used a simple SentenceBERT embedding model all-MiniLM-L6-v2 \parencite{reimers2019sentence} to vectorise captions and claims. 
\begin{table}[ht]
\center
\begin{tabular}{lc}
\hline
Statements & Cosine similarity between caption and statement\\
\hline
Entailed  & 0.3337\\
Refuted & 3320\\
\hline
\end{tabular}
\caption{Cosine similarity between table captions and statements.}
\label{tab:tabfact_caption_cosine_sim}
\end{table}

\section{Baseline Fact-Verification on TabFact Results}
The results of the simple table fact-verification baseline pipelines are shown in Table~\ref{tab:tabfact_baseline_results}. As you can see from the table, most of the baselines achieved modest performance, not much higher than chance. Logistic regression is generally the poorest, highlighting the need for a model which can make sense of the structure. The BERT encoder model produced superior metrics, with accuracy near 60 and F1 in the 60s. It shows that the contextual representation of the table and claim contributed to success. It actually produced the highest recall of any baseline, but at the expense of many false positives, as precision is low. The ordering of table and claim had little difference; the F1 and accuracy are nearly identical. Performance would have been affected by 18\% of its data being truncated, as it only allows up to 512 tokens. The decoder model GPT-2 has more parameters and allows 1024 tokens, so only about 4\% were truncated. However, it did not outperform BERT. This could be due to not optimising hyperparameters enough; the literature suggests its potential is higher, but clearly both do not look too promising at achieving strong fact verification performance with a complete table and claim concatenation. 

\begin{table}[ht]
\centering
\small
\begin{tabular}{p{5.5cm} c c c c}
\hline
\textbf{Fact Verification Baseline} & \textbf{Test Accuracy} & \textbf{Test F1} & \textbf{Test Precision} & \textbf{Test Recall} \\
\hline
Logistic regression T+C & 0.5227 & 0.6253 & 0.5165 & 0.7922 \\
Logistic regression C+T & 0.5229 & 0.6263 & 0.5169 & 0.7952 \\[4pt]
BERT base T+C & 0.5768 & 0.6711 & 0.5508 & \textbf{0.8588} \\
BERT base C+T & 0.5888 & 0.6715 & 0.5611 & 0.8358 \\[4pt]
GPT-2 T+C & 0.5604 & 0.6518 & 0.5416 & 0.8185 \\
GPT-2 C+T & 0.5613 & 0.6253 & 0.5480 & 0.7279 \\[4pt]
Qwen3-0.6B zero-shot T+C & 0.5561 & 0.5825 & 0.5526 & 0.6159 \\
Qwen3-0.6B zero-shot C+T & 0.5317 & 0.5420 & 0.5331 & 0.5513 \\
Qwen3-0.6B one-shot T+C & 0.5534 & 0.5183 & 0.5661 & 0.4780 \\
Qwen3-0.6B one-shot C+T & 0.5227 & 0.5258 & 0.5258 & 0.5150 \\
Qwen3-0.6B two-shot T+C & 0.5466 & 0.5506 & 0.5488 & 0.5523 \\
Qwen3-0.6B two-shot C+T & 0.5187 & 0.5679 & 0.5175 & 0.6291 \\[4pt]
Qwen3-0.6B zero-shot with thinking (sample of 426) T+C & 0.6556 & 0.6260 & \textbf{0.6848} & 0.5765 \\
Qwen3-0.6B zero-shot with thinking (sample of 426) C+T & 0.6650 & 0.6562 & 0.6477 & 0.6649 \\[4pt]
Qwen3-0.6B with LoRA PEFT T+C & \textbf{0.7197} & \textbf{0.7511} & 0.6784 & 0.8414 \\
Qwen3-0.6B with LoRA PEFT C+T & 0.6088 & 0.6125 & 0.6101 & 0.6149 \\
\hline
\end{tabular}
\caption{Fact verification baseline results across models and input orderings (\texttt{T+C} vs.\ \texttt{C+T}).}
\label{tab:tabfact_baseline_results}
\end{table}

\vspace{\baselineskip}
\noindent The Qwen few-shot results show that adding additional table and claim examples was not very effective. There is a good chance with more effort put into hyperparameter selection, primarily the selection of the entailed and refuted evidence as I did some brief searching for table representations on small samples that would produce a result superior to random guessing before I would use it on the test set. Before I did this, I ran one iteration of one-shot, which produced a result of 49\%. So it is clear that selection is important. The more shots added, the worse the accuracy was, and Qwen produced superior results for table-before-claim. Two shots did give the best F1 for C+T, though. 

\vspace{\baselineskip}
\noindent The zero-shot test with thinking showed quite impressive results, very similar to both configurations of the input, showing accuracy, precision, and F1 in the mid-60s. Claim after table had the highest precision out of any baseline, but table after claim had a superior F1. This can not be trusted too much, though, as it was with a small sample; it was randomised and stratified, so I would say it is still a promising result. 

\vspace{\baselineskip}
\noindent The stand-out baselines were the LoRa fine-tuned Qwen models. It produced decent results for claim-then-table, with all metrics being in the low 60s. The result with the claim concatenated after the table, though, did very well: accuracy and F1 in the low and mid 70s, respectively, with moderate precision but impressive recall, suggesting some false positives. It is notable that this was achieved with only a single epoch as it took very long to fine-tune, and with 16\% of the samples truncated. With multiple epoch runs, and a row-evidence-sourced table, the ceiling for its performance could be higher than what we see here. It is clear that table-then-claim representation, which TabFact also found to be the best, with Qwen fine-tuned with LoRa, will certainly have to be examined for possible more complex pipelines later as it shows the most promise. 
